% s00-front.tex —— 前置：摘要与译本说明
% 摘要对应原文 p.1 的 Abstract；原文日期与资助信息见 p.1 脚注

\section*{摘要}
\addcontentsline{toc}{section}{摘要}

本文概述作者所提出的局部 Langlands 对应几何化猜想。\yizhuw{%
原文摘要仅一句：\emph{This article is an overview of the geometrization conjecture for the
local Langlands correspondence formulated by the author.} 原刊日期为 2016 年 2 月 2 日；
作者受法国国家科研署项目 ANR-14-CE25 ``PerCoLaTor'' 资助。%
}

\section*{译本说明}
\addcontentsline{toc}{section}{译本说明}

\subsection*{一、底本}

\begin{itemize}
\item Laurent Fargues, \emph{Geometrization of the local Langlands correspondence: an overview},
2016 年 2 月 2 日，46 页。
\item 底本为作者自排的 pdf\LaTeX{} 文稿，无出版方版式；印刷页码与 PDF 页码一致（第 1 页即 PDF 第 1 页）。
\item 译文正文中的边注 \textsf{p.$N$} 标出该段落在底本中的起始页，便于逐页核对。
\end{itemize}

\subsection*{二、体例}

\begin{itemize}
\item 章节、定理、命题、定义、注、猜想的编号与底本完全一致：底本按节连续编号
（如定义 1.5、命题 1.6、注 1.7 依次相连），译本沿用同一计数器，故编号可直接对照。
\item 公式编号亦与底本一致；底本无编号的展示公式，译本同样不编号。
\item 译者注以脚注形式排于页脚，楷体，并标注来源等级：
A 为原文或作者本人论文、经典教材、arXiv 预印本；B 为 nLab、Stacks Project、ProofWiki 等白名单站点
（须与 A 类交叉核实）；I 为本译丛既有译注、术语表与资料卡。无外部依据者标「译者按（无外部依据）」。
\item 底本明显的笔误，译本一律照排原文并加注说明，不径改正文；全部条目登记于
\texttt{工程文件/work/疑点清单.tsv}。
\item 术语依 \texttt{工程文件/work/glossary.tsv} 冻结；未定名者保留原文并加注。
\end{itemize}

\subsection*{三、术语约定}

本译本沿用本译丛在 p-进 Hodge 理论与几何 Langlands 方向已冻结的译名，主要有：
diamond 译「菱形」，perfectoid 译「完满」，untilt 与 pro-étale 保留原文，
$\Bun_G$ 与 $\Gr$ 保留记号，Hecke eigensheaf 译「Hecke 特征层」，
perverse sheaf 译「反常层」，$\sigma$-conjugacy 译「$\sigma$-共轭」。
完整对照见术语表；同一术语在全书内不得改译。

\subsection*{四、分工与责任}

本译本由数译工程组织完成：底本逐页切分后分章精读与翻译，再由统一体例合成、逐页核对公式，
并逐条登记疑点。译文对底本的忠实度以「不改原文、只加批注」为限；
读者若发现译文与底本图像不符，请以底本为准并告知我们修订。
